% GENERATED FILE. Edit canonical lessons, then run npm run generate:books.
% canonical-source-hash: fnv1a64:0b0f3792ff539df8
% canonical-lessons: GU-C374-avrodh, GU-C374-vighna, GU-C374-musibat, GU-C374-khalvu, GU-C374-nivarvu

\chapter{Obstacles}
\label{ch:gu-obstacles}
\hlchaptermodality{374}

\begin{chapteropening}
\textbf{By the end of this chapter:} \emph{I can say what might stand in the way and how to stop it.}

Complete the last lesson of chapter 374: I can say what might stand in the way and how to stop it.
\end{chapteropening}

\section[\texorpdfstring{\gu{અવરોધ} (avrodh)}{અવરોધ (avrodh)}]{\gu{અવરોધ} (avrodh) — an obstruction}

\label{lesson:GU-C374-avrodh}



\begin{quote}
Before the new one: say the Gujarati for a difficulty, then the Gujarati for to get in the way.
\end{quote}

\subsection*{You'll want to know: \gu{અવરોધ}}
\textbf{\gu{અવરોધ}} — \emph{avrodh} — \textquotedblleft{}an obstruction\textquotedblright{}.

\begin{quote}
\gu{જો} \gu{રસ્તામાં} \gu{કોઈ} \gu{અવરોધ} \gu{હશે} \gu{તો} \gu{આપણે} \gu{બીજા} \gu{રસ્તે} \gu{જઈશું}.
\end{quote}

\begin{quote}
\emph{jo rastāmā̃ koī avrodh haśe to āpaṇe bījā raste jaīśũ.}
\end{quote}

\begin{quote}
If there is something blocking the road, we will take another way.
\end{quote}

A masculine noun.

\begin{cousinweb}
\textbf{\gu{અવરોધ}} (\emph{avrodh}) is a tatsama, borrowed from Sanskrit \emph{avarodha}, \textquotedblleft{}hindrance, blockade\textquotedblright{} (Monier-Williams).
\end{cousinweb}

\begin{grammarlens}[title={What you've built}]
One more word for saying what happens if.
\end{grammarlens}

\subsection*{Your turn}
\begin{itemize}
\raggedright
  \item \emph{Say it:} \emph{avrodh}
  \item \emph{Say it:} \emph{avrodh}, once more
  \item \emph{Say it:} \emph{jo rastāmā̃ koī avrodh haśe to āpaṇe bījā raste jaīśũ.} — If there is something blocking the road, we will take another way.
  \item \emph{Recall:} say the Gujarati for a difficulty, then the Gujarati for to get in the way, then say \emph{avrodh} again
\end{itemize}

\begin{tcolorbox}[breakable,colback=teal!4,colframe=teal!35!black,title={Before you move on}]
What does \gu{અવરોધ} mean? (\textquotedblleft{}An obstruction\textquotedblright{}.) Say it once more.
\end{tcolorbox}
\section[\texorpdfstring{\gu{વિઘ્ન} (vighna)}{વિઘ્ન (vighna)}]{\gu{વિઘ્ન} (vighna) — a hindrance}

\label{lesson:GU-C374-vighna}



\begin{quote}
Before the new one: say the Gujarati for to get in the way, then the Gujarati for an obstruction.
\end{quote}

\subsection*{You'll want to know: \gu{વિઘ્ન}}
\textbf{\gu{વિઘ્ન}} — \emph{vighna} — \textquotedblleft{}a hindrance\textquotedblright{}.

\begin{quote}
\gu{જો} \gu{કોઈ} \gu{વિઘ્ન} \gu{આવશે} \gu{તો} \gu{અમે} \gu{લગ્ન} \gu{મુલતવી} \gu{રાખીશું}.
\end{quote}

\begin{quote}
\emph{jo koī vighna āvśe to ame lagna mulatvī rākhīśũ.}
\end{quote}

\begin{quote}
If some hitch comes up, we will put off the wedding.
\end{quote}

A neuter noun.

\begin{cousinweb}
\textbf{\gu{વિઘ્ન}} (\emph{vighna}) is a tatsama, borrowed from Sanskrit \emph{vighna}, \textquotedblleft{}a breaker; an obstacle\textquotedblright{} (Monier-Williams).
\end{cousinweb}

\begin{grammarlens}[title={What you've built}]
One more word for saying what happens if.
\end{grammarlens}

\subsection*{Your turn}
\begin{itemize}
\raggedright
  \item \emph{Say it:} \emph{vighna}
  \item \emph{Say it:} \emph{vighna}, once more
  \item \emph{Say it:} \emph{jo koī vighna āvśe to ame lagna mulatvī rākhīśũ.} — If some hitch comes up, we will put off the wedding.
  \item \emph{Recall:} say the Gujarati for to get in the way, then the Gujarati for an obstruction, then say \emph{vighna} again
\end{itemize}

\begin{tcolorbox}[breakable,colback=teal!4,colframe=teal!35!black,title={Before you move on}]
What does \gu{વિઘ્ન} mean? (\textquotedblleft{}A hindrance\textquotedblright{}.) Say it once more.
\end{tcolorbox}
\section[\texorpdfstring{\gu{મુસીબત} (musībat)}{મુસીબત (musībat)}]{\gu{મુસીબત} (musībat) — trouble}

\label{lesson:GU-C374-musibat}



\begin{quote}
Before the new one: say the Gujarati for an obstruction, then the Gujarati for a hindrance.
\end{quote}

\subsection*{You'll want to know: \gu{મુસીબત}}
\textbf{\gu{મુસીબત}} — \emph{musībat} — \textquotedblleft{}trouble\textquotedblright{}.

\begin{quote}
\gu{જો} \gu{તમે} \gu{મોડું} \gu{કરશો} \gu{તો} \gu{મુસીબતમાં} \gu{પડશો}.
\end{quote}

\begin{quote}
\emph{jo tame moḍũ karśo to musībatmā̃ paḍśo.}
\end{quote}

\begin{quote}
If you are late, you will land in trouble.
\end{quote}

A feminine noun.

\begin{cousinweb}
\textbf{\gu{મુસીબત}} (\emph{musībat}) is usually traced to Arabic \emph{muṣība}, \textquotedblleft{}a misfortune\textquotedblright{}, through Persian.
\end{cousinweb}

\begin{grammarlens}[title={What you've built}]
One more word for saying what happens if.
\end{grammarlens}

\subsection*{Your turn}
\begin{itemize}
\raggedright
  \item \emph{Say it:} \emph{musībat}
  \item \emph{Say it:} \emph{musībat}, once more
  \item \emph{Say it:} \emph{jo tame moḍũ karśo to musībatmā̃ paḍśo.} — If you are late, you will land in trouble.
  \item \emph{Recall:} say the Gujarati for an obstruction, then the Gujarati for a hindrance, then say \emph{musībat} again
\end{itemize}

\begin{tcolorbox}[breakable,colback=teal!4,colframe=teal!35!black,title={Before you move on}]
What does \gu{મુસીબત} mean? (\textquotedblleft{}Trouble\textquotedblright{}.) Say it once more.
\end{tcolorbox}
\section[\texorpdfstring{\gu{ખાળવું} (khāḷvũ)}{ખાળવું (khāḷvũ)}]{\gu{ખાળવું} (khāḷvũ) — to keep in check}

\label{lesson:GU-C374-khalvu}



\begin{quote}
Before the new one: say the Gujarati for a hindrance, then the Gujarati for trouble.
\end{quote}

\subsection*{You'll want to know: \gu{ખાળવું}}
\textbf{\gu{ખાળવું}} — \emph{khāḷvũ} — \textquotedblleft{}to keep in check\textquotedblright{}.

\begin{quote}
\gu{જો} \gu{આપણે} \gu{નદીનું} \gu{પાણી} \gu{ખાળીશું} \gu{તો} \gu{ગામ} \gu{બચી} \gu{જશે}.
\end{quote}

\begin{quote}
\emph{jo āpaṇe nadīnũ pāṇī khāḷīśũ to gām bacī jaśe.}
\end{quote}

\begin{quote}
If we hold back the river water, the village will be saved.
\end{quote}

Past, with \gu{તેણે} (he or she …): \gu{તેણે} \gu{આંસુ} \gu{ખાળ્યાં}, \emph{teṇe ā̃su khāḷyā̃}.

\begin{cousinweb}
\textbf{\gu{ખાળવું}} (\emph{khāḷvũ}) is inherited from Sanskrit \emph{skhālayati}, \textquotedblleft{}causes to falter, stops\textquotedblright{}; CDIAL lists the Gujarati verb, \textquotedblleft{}to hinder, to stop\textquotedblright{} (CDIAL 13664).
\end{cousinweb}

\begin{grammarlens}[title={What you've built}]
One more word for saying what happens if.
\end{grammarlens}

\subsection*{Your turn}
\begin{itemize}
\raggedright
  \item \emph{Say it:} \emph{khāḷvũ}
  \item \emph{Say it:} \emph{khāḷvũ}, once more
  \item \emph{Say it:} \emph{jo āpaṇe nadīnũ pāṇī khāḷīśũ to gām bacī jaśe.} — If we hold back the river water, the village will be saved.
  \item \emph{Recall:} say the Gujarati for a hindrance, then the Gujarati for trouble, then say \emph{khāḷvũ} again
\end{itemize}

\begin{tcolorbox}[breakable,colback=teal!4,colframe=teal!35!black,title={Before you move on}]
What does \gu{ખાળવું} mean? (\textquotedblleft{}To keep in check\textquotedblright{}.) Say it once more.
\end{tcolorbox}
\section[\texorpdfstring{\gu{નિવારવું} (nivārvũ)}{નિવારવું (nivārvũ)}]{\gu{નિવારવું} (nivārvũ) — to avert}

\label{lesson:GU-C374-nivarvu}



\begin{quote}
Before the new one: say the Gujarati for trouble, then the Gujarati for to keep in check.
\end{quote}

\subsection*{You'll want to know: \gu{નિવારવું}}
\textbf{\gu{નિવારવું}} — \emph{nivārvũ} — \textquotedblleft{}to avert\textquotedblright{}.

\begin{quote}
\gu{જો} \gu{આપણે} \gu{હાથ} \gu{ધોઈશું} \gu{તો} \gu{ઘણા} \gu{રોગ} \gu{નિવારી} \gu{શકીશું}.
\end{quote}

\begin{quote}
\emph{jo āpaṇe hāth dhoīśũ to ghaṇā rog nivārī śakīśũ.}
\end{quote}

\begin{quote}
If we wash our hands, we can prevent many illnesses.
\end{quote}

Past, with \gu{તેણે} (he or she …): \gu{તેણે} \gu{મોટી} \gu{આપત્તિ} \gu{નિવારી}, \emph{teṇe moṭī āpatti nivārī}.

\begin{cousinweb}
\textbf{\gu{નિવારવું}} (\emph{nivārvũ}) is inherited from Sanskrit \emph{nivārayati}, \textquotedblleft{}wards off, restrains\textquotedblright{}; CDIAL lists the Gujarati verb, \textquotedblleft{}to turn back, to avert\textquotedblright{} (CDIAL 7419).
\end{cousinweb}

\begin{grammarlens}[title={What you've built}]
One more word for saying what happens if.
\end{grammarlens}

\subsection*{Your turn}
\begin{itemize}
\raggedright
  \item \emph{Say it:} \emph{nivārvũ}
  \item \emph{Say it:} \emph{nivārvũ}, once more
  \item \emph{Say it:} \emph{jo āpaṇe hāth dhoīśũ to ghaṇā rog nivārī śakīśũ.} — If we wash our hands, we can prevent many illnesses.
  \item \emph{Recall:} say the Gujarati for trouble, then the Gujarati for to keep in check, then say \emph{nivārvũ} again
\end{itemize}

\begin{tcolorbox}[breakable,colback=teal!4,colframe=teal!35!black,title={Before you move on}]
What does \gu{નિવારવું} mean? (\textquotedblleft{}To avert\textquotedblright{}.) Say it once more.
\end{tcolorbox}
